\section{Sikkerhetsarkitektur}
\label{sec:design-sikkerhet}
Dette delkapittelet beskriver sikkerhetsdesignet i VideoAPI. Først gjennomgås autentisering og tilgangskontroll, inkludert valg av Keycloak og OAuth2-flyt. Deretter beskrives claim-struktur og tokenvalidering, før det avsluttes med en gjennomgang av utvalgte angrepsscenarier og tilhørende kontroller. En STRIDE-basert trusselmodell finnes i delkapittel~\ref{sec:trusselmodell}.

\subsection{Autentisering og tilgangskontroll}

Autentisering i VideoAPI er basert på OAuth2-standarden, noe som var et eksplisitt krav fra oppdragsgiver. Valget av OAuth2 gir en industristandard tilnærming til tilgangskontroll som er vidt støttet og godt dokumentert, og som muliggjør fremtidig integrasjon med andre systemer uten behov for endringer i API-ets grensesnitt.

\subsection{Valg av Keycloak som identitetsleverandør}

Som OAuth2-leverandør ble Keycloak valgt. Bakgrunnen for dette valget er at HDO allerede benytter Keycloak som identity provider i sin eksisterende infrastruktur. Ved å ta i bruk det samme systemet unngikk gruppen å innføre ny teknologi i driftsmiljøet, og oppdragsgiver kan selv ta over og tilpasse konfigurasjonen i etterkant uten å måtte sette seg inn i et ukjent verktøy.

Autentiseringsmekanismen er implementert ved bruk av OAuth 2.0-flyten \textit{Client Credentials Grant}, også omtalt som \gls{m2m}-autentisering. Denne flyten er passende for et API uten sluttbrukerpålogging, der klientsystemer autentiserer seg med en klient-ID og klienthemmelighet for å hente et JWT-token. Tokenet brukes deretter som \texttt{Bearer}-token i alle påfølgende API-kall.

\subsection{Minimal claim-struktur}

Et bevisst designvalg i autentiseringsarkitekturen er å holde claim-strukturen enkel. Det er kun én claim som kreves for tilgang: \texttt{api\_access = true}. Det er ikke benyttet OAuth2-scopes.

Denne forenklingen er begrunnet i prosjektets kontekst: autentiseringsoppsettet er ment som et midlertidig rammeverk som oppdragsgiver selv skal overta og konfigurere videre med sin egen Keycloak-instans. Å designe et komplekst scope-hierarki som senere ville blitt erstattet ville ført til unødvendig arbeid uten varig verdi. Den enkelte \texttt{api\_access}-claimen fungerer derfor som en enkel binær tilgangskontroll, hvor en klient enten gis tilgang eller avvises.

Keycloak-klienter som skal ha tilgang til API-et må konfigureres med en protokollmapper som legger til denne claimet i tokens som utstedes. Selve policy-sjekken skjer globalt i API-ets autorisasjonsmiddleware, slik at alle beskyttede endepunkter automatisk krever gyldig token med riktig claim uten at dette må defineres eksplisitt per endepunkt.

\subsection{Validering og sikkerhetsnivå}

JWT-tokens valideres mot Keycloaks JWKS-endepunkt, som betyr at signaturen verifiseres med Keycloaks offentlige nøkkel ved hvert kall. Audience- og issuer-validering er i det nåværende oppsettet deaktivert, da disse parameterne vil variere avhengig av hvilken Keycloak-instans oppdragsgiver velger å benytte. Signaturvalidering sikrer at tokens ikke kan forfalskes, og at kun tokens utstedt av den konfigurerte Keycloak-realmen godtas.

\subsection{Utvalgte angrep og kontroller}

Token-endepunktet \texttt{POST /auth/token} er beskyttet med en skjerpet
rate-limit på 10 forespørsler per minutt per IP-adresse, sammenlignet med
den generelle grensen på 60 forespørsler per minutt. Dette er et tiltak
mot credential-stuffing og ressursuttømming mot autentiseringslaget.

% Autorisasjonsmodellen er flat: enhver klient med gyldig JWT og
% \texttt{api\_access}-claim har samme rettigheter. Rom-ID-er er derfor ikke
% personlige hemmeligheter i seg selv. Dersom en klient kjenner en annen
% klients rom-ID, kan den i prinsippet kalle rom-endepunktene, med mindre
% dette hindres utenfor VideoAPI. Rom-spesifikk isolasjon mellom
% integrasjonsparter er utenfor proof-of-conceptets omfang, men dokumenteres
% her som en bevisst avgrensning.

Autorisasjonsmodellen er flat: enhver klient med gyldig JWT og \texttt{api\_access}-claim tildeles samme rettigheter. Rom-ID-er betraktes derfor ikke som personlige hemmeligheter i seg selv. Dersom en klient kjenner en annen klients rom-ID, kan den i prinsippet benytte rom-endepunktene, med mindre dette begrenses av mekanismer utenfor VideoAPI. Romspesifikk isolasjon mellom integrasjonsparter ligger utenfor POC-ets omfang og er dokumentert som en bevisst avgrensning.

En STRIDE-basert oppsummering av trusler og tiltak finnes i
delkapittel~\ref{sec:trusselmodell}.